\documentclass{article}
\usepackage{amsmath,hyperref}
\begin{document}
\title{Resuming a continuous optical receipt stream}
\maketitle

\paragraph{Result.}
An 11-bit optical UART stream can carry an 84-symbol self-synchronizing
alphabet without a new delimiter or a longer word. This is the largest such
alphabet within the published start, parity and stop format. The result assumes
one trusted emitter, a contiguous stream of valid words, and correct bits at a
known bit cadence, but no known word phase after obscuration. It says when the
receiver can recognize a complete word. It does not certify the emitter, the
freshness of the receipt, or a physical handoff.

\paragraph{Source and baseline.}
Wang, Ng, Henderson and Mahony
(\href{https://arxiv.org/abs/2208.01710}{IROS 2022}, sections 3.1,
3.2.4--5 and 4.1) transmit a zero start bit, seven payload bits,
one even parity bit and two one stop bits. They look for a falling edge near
the expected next word boundary. Incorrectly positioned edges and parity tests
are discussed probabilistically; the reported outdoor test transmits a sonnet
twice, separated by a long high interval. The
\href{https://github.com/ziweiWWANG/Event-Beacon-Communication/blob/00be0f34eea83d2a3508706a4d0a922c4047062e/decoder.m}{released decoder}
starts from a long-high candidate. The midstream phase model below may have
no such interval after the view reopens. The result is not a rerun of their
physical experiments or a failure of their start-of-message benchmark.

\paragraph{Obstruction.}
Write \(F(x)\) for the 11-bit word encoding payload \(x\). In a continued
repetition of \(F(83)=01010011011\), the falling edge at offset eight can
be treated as the start of \(F(106)=01101010011\) on every repetition. Even a
strict receiver that checks the start and stop bits as well as parity accepts that second
parse. If the receiver sees only the indefinitely repeated bits and has no
word-phase side information, the two parses are observationally identical.
The same fixed offset is a falling edge each cycle, so waiting for more
such cycles does not make a falling-edge search slip. The replay finds 96
directed aliases across 84 of the 128 payloads. These are algebraic counts,
not frequencies in the authors' data.

A choice that fixes repetitions alone need not fix changes of value. The
splice of \(F(32)\) and \(F(39)\) at offset seven gives \(F(114)\); those three
words are from different rotation classes. Choosing the smallest payload in
every class likewise leaves 416 spurious ordered splices. Choosing the
largest leaves 376. An alphabet with no internal \(110\) marker is a
simple sufficient restriction: the marker can then occur only across a
real boundary (two stop ones and the next start zero). That choice yields
44 symbols in this format.

\paragraph{Maximum alphabet.}
For an admissible alphabet \(C\) and \(a,b\) in \(C\), let
\(W_s(a,b)=a[s:11]\,b[0:s]\) with zero-based offsets. The exact condition
is \(W_s(a,b)\notin C\) for \(1\leq s\leq 10\). Necessity uses the specified
receiver, which accepts a window whenever it is a codeword. Sufficiency
holds for arbitrary sequences because any wrong-phase 11-bit window
intersects exactly two successive words. This is the classical
\href{https://doi.org/10.4153/CJM-1958-023-9}{comma-free code} criterion
applied to the smaller set permitted by this UART format.

In particular two rotations cannot both occur in \(C\) (use \(a=b\)).
All 11-bit UART words are nonconstant and therefore primitive since 11
is prime. There are 84 rotation classes meeting the format, as counted by
choosing the least of the eleven rotations of each of the 128 explicit
words. The accompanying list picks one per class and forbids every
wrong-phase splice, verified by the accompanying complete enumeration.
The bound thus does not depend on trusting the optimization used to find
the list. It applies within the unchanged UART format and to this
unconstrained sequence of payloads and the specified receiver. With a
framed preamble, previously known phase, or a grammar restricting
successive values, this upper bound is not claimed.

\paragraph{Usable evidence cut.}
A sliding-window receiver can recognize the first wholly visible word
in at most 21 observed bits, from any integer bit cut, provided the
stream continues. A cut at offset zero requires 11 bits; a cut at
\(s>0\) requires \(22-s\). This is tight for that receiver and that
output; no claim is made that another receiver cannot infer a restricted
message earlier. The bound counts sampled bits only. It excludes fault
detection, readiness of the emitter, delivery of events to the decoder,
subsequent processing, and any need to repeat a particular receipt.
A freshly visible stream can still contain an obsolete receipt. A
consumer must separately bind the decoded symbol to the current issuer,
duty, generation and effect. The phase decoder supplies none of these
claims solely by accepting a word.

Every aligned single-bit substitution fails exact alphabet membership,
because all words have even weight. That protection does not compose
with an unknown phase: the offset-eight window from words 0 and 1 is
\(01100000001\); one bit flip makes it \(01100000011=F(96)\), a
selected word. Bit recovery, occlusion length, association of events
with one emitter, and corruption or insertion/deletion remain outside the
guarantee. A preamble or a code with a tested distance from all
wrong-phase splices is a viable alternative if some errors must be
tolerated. Neither this calculation nor a nominal baud rate establishes
the timely end-to-end observation required by the relay test.

\paragraph{Verification and project use.}
A textual receiver was checked against a packed-integer window scan
from all 77,616 cuts in all pairs of selected words. The replay checks
all 70,560 wrong-phase splices, all 924 aligned single flips, the
class count and the cited counterexamples. Lean checks the three
explicit arithmetic witnesses and the observation count, not the full
codebook search. This provides an R1 read-cut condition for one optical
framing design, not the independently selected non-radio replacement of
an authoritative receipt.
\end{document}
